\section{Discussion}

\subsection{What Mode 3 Tells Us}

The 24\% Mode~3 finding reveals a substantial blind spot in RLHF training signal. In one quarter of preference battles, reward models express high confidence on samples where humans disagree. This systematic overconfidence may propagate through RLHF training, reinforcing model behaviors that do not reflect genuine human consensus.

\subsection{What Mode 3 Does Not Tell Us}

We have not identified \emph{why} overconfident misalignment occurs. Semantic divergence is falsified (Mode~3 pairs are \emph{more} similar, not less). Prompt subjectivity is unsupported (Mode~3 pervades all task types with negligible enrichment in subjective prompts).

\subsection{Competing Explanations}

Several alternative mechanisms merit investigation:
\begin{itemize}
    \item \textbf{Style/tone differences}: Embeddings may miss presentation-level distinctions that humans weight heavily
    \item \textbf{RM feature sensitivity}: Reward models may attend to features that humans largely ignore
    \item \textbf{Embedding limitations}: Larger embedding models may reveal semantic differences our analysis missed
\end{itemize}

\subsection{Limitations}

\begin{itemize}
    \item Median-split thresholding: Alternative clustering methods may shift mode boundaries
    \item Near-uniform distribution may partly reflect methodological choice rather than ground truth
    \item Single RM variance proxy: Full ensemble validation pending
    \item Chatbot Arena specific: Generalization to other preference datasets unknown
    \item Correlational only: No causal claims
\end{itemize}
